% Part: first-order-logic
% Chapter: tableaux
% Section: quantifier-rules

\documentclass[../../../include/open-logic-section]{subfiles}

\begin{document}

\olfileid{fol}{tab}{qrl}

\olsection{পরিমাণসূচকের বিধি}

\subsection{$\lforall$-এর বিধি}

\begin{defish}
\AxiomC{\sFmla{\True}{\lforall[x][!A(x)]}}
\RightLabel{\TRule{\True}{\lforall}}
\UnaryInfC{\sFmla{\True}{!A(t)}}
\DisplayProof
\hfill
\AxiomC{\sFmla{\False}{\lforall[x][!A(x)]}}
\RightLabel{\TRule{\False}{\lforall}}
\UnaryInfC{\sFmla{\False}{!A(a)}}
\DisplayProof
\end{defish}

\TRule{\True}{\lforall}-এ $t$ একটি বদ্ধ পদ (অর্থাৎ চলরাশিবিহীন
পদ)। \TRule{\False}{\lforall}-এ $a$ এমন একটি !!a{constant}, যা
\TRule{\False}{\lforall} বিধির ওপরের শাখার কোথাও থাকতে পারবে না।
$a$-কে \TRule{\False}{\lforall} অনুমানের \emph{আইগেনচল} বলি।%
\footnote{ওপরের বিধিতে $a$ আসলে একটি !!a{constant} হলেও ঐতিহাসিক
কারণে “আইগেনচল” শব্দটি ব্যবহার করি।}
% BN-SRC-044 TARGET-CORRECTION-BEGIN
\par\smallskip\noindent\textnormal{\small\textbf{উৎস-সংশোধন (BN-SRC-044):} হিমায়িত ইংরেজি বিধি ও পরের উল্লেখে সার্বিক পরিমাণসূচকের কনফিগারযোগ্য ম্যাক্রো \texttt{\string\lforall}-এর বদলে দুবার কাঁচা \texttt{\string\forall} দেওয়া হয়েছে। অধ্যায়ের অন্য সব বিধি-লেবেলের সঙ্গে মিলিয়ে বাংলা লক্ষ্যে \texttt{\string\lforall} পুনরুদ্ধার করা হয়েছে।} % BN-SRC-044 TARGET-CORRECTION-END

\subsection{$\lexists$-এর বিধি}

\begin{defish}
\AxiomC{\sFmla{\True}{\lexists[x][!A(x)]}}
\RightLabel{\TRule{\True}{\lexists}}
\UnaryInfC{\sFmla{\True}{!A(a)}}
\DisplayProof
\hfill
\AxiomC{\sFmla{\False}{\lexists[x][!A(x)]}}
\RightLabel{\TRule{\False}{\lexists}}
\UnaryInfC{\sFmla{\False}{!A(t)}}
\DisplayProof
\end{defish}

এখানেও $t$ একটি বদ্ধ পদ এবং $a$ এমন একটি !!a{constant}, যা
\TRule{\True}{\lexists} বিধির ওপরের শাখায় থাকে না। $a$-কে
\TRule{\True}{\lexists} অনুমানের \emph{আইগেনচল} বলি।

\TRule{\False}{\lforall} বা \TRule{\True}{\lexists} অনুমানের ওপরের
শাখায় আইগেনচল না থাকার শর্তটিকে \emph{আইগেনচল-শর্ত} বলা হয়।

\begin{explain}
মনে করো, $!A$ যদি এমন একটি !!a{formula} হয় যাতে !!{variable}~$x$
মুক্ত, তবে প্রচলিত রীতি অনুযায়ী তা~$!A(x)$ লিখি। একই প্রসঙ্গে
$!A(t)$ হলো~$\Subst{!A}{t}{x}$-এর সংক্ষিপ্ত রূপ। তাই
\TRule{\False}{\lexists} বিধিটি এভাবেও লেখা যায়:
\begin{prooftree}
  \AxiomC{\sFmla{\False}{\lexists[x][!A]}}
  \RightLabel{\TRule{\False}{\lexists}}
  \UnaryInfC{\sFmla{\False}{\Subst{!A}{t}{x}}}
\end{prooftree}
লক্ষ করো, $t$ আগে থেকেই~$!A$-তে থাকতে পারে; যেমন, $!A$ হতে পারে
$\Atom{\Obj P}{t,x}$। তাই
$\sFmla{\False}{\lexists[x][\Atom{\Obj P}{t,x}]}$ থেকে
$\sFmla{\False}{\Atom{\Obj P}{t,t}}$ অনুমান করা
\TRule{\False}{\lexists}-এর সঠিক প্রয়োগ। কিন্তু
\TRule{\False}{\lforall} ও~\TRule{\True}{\lexists}-এর
আইগেনচল-শর্তে !!{constant}~$a$-কে $!A$-তে না থাকতে হয়। অতএব
$\sFmla{\False}{\lforall[x][\Atom{\Obj P}{a,x}]}$ থেকে
$\sFmla{\False}{\Atom{\Obj P}{a,a}}$
ব্যবহার করে~$\TRule{\False}{\lforall}$ সঠিকভাবে অনুমান করা যায় না।
\end{explain}

\begin{explain}
\TRule{\True}{\lforall} ও \TRule{\False}{\lexists}-এ $t$-কে বদ্ধ
পদ হতে হয়, তবে তার ওপর আর কোনো নতুনত্ব-শর্ত নেই। অন্যদিকে
\TRule{\True}{\lexists} ও \TRule{\False}{\lforall} বিধিতে
আইগেনচল-শর্ত অনুসারে !!{constant}~$a$ সংশ্লিষ্ট অনুমানের ওপরের কোনো
শাখায় থাকতে পারে না। পদ্ধতির বিশুদ্ধতা নিশ্চিত করতে এই শর্ত জরুরি।
শর্তটি না থাকলে
$\lexists[x][\formula{A}(x)] \lif \lforall[x][\formula{A}(x)]$-এর
জন্য নিচের বদ্ধ !!{tableau}-টি পাওয়া যেত:
\begin{center}
\begin{tableau}{}
  [\sFmla{\False}{\lexists[x][\formula{A}(x)] \lif \lforall[x][\formula{A}(x)]}, just=\TAss
    [\sFmla{\True}{\lexists[x][\formula{A}(x)]},
      just={\TRule{\False}{\lif}[1]}
      [\sFmla{\False}{\lforall[x][\formula{A}(x)]},
        just={\TRule{\False}{\lif}[1]}
        [\sFmla{\True}{\formula{A}(a)},
          just={\TRule{\True}{\lexists}[2]}
          [\sFmla{\False}{\formula{A}(a)},
            just={\TRule{\False}{\lforall}[3]}, close]
        ]
      ]
    ]
  ]
\end{tableau}
\end{center}
কিন্তু $\lexists[x][\formula{A}(x)] \lif
\lforall[x][\formula{A}(x)]$ বৈধ নয়।
% BN-SRC-045 TARGET-CORRECTION-BEGIN
\par\smallskip\noindent\textnormal{\small\textbf{উৎস-স্পষ্টীকরণ (BN-SRC-045):} হিমায়িত ইংরেজি গদ্যের শেষ ব্যাখ্যায় $t$-র ওপর “কোনো বিধিনিষেধ নেই” বলা হলেও আগেই $t$-কে বদ্ধ পদ বলা হয়েছে। বাংলা গদ্যে তাই সঙ্গত অর্থটি—বদ্ধতার শর্ত ছাড়া কোনো অতিরিক্ত নতুনত্ব-শর্ত নেই—স্পষ্ট করা হয়েছে।} % BN-SRC-045 TARGET-CORRECTION-END
\end{explain}

\end{document}
